\documentclass[11pt,a4paper]{article}
\usepackage[margin=25mm]{geometry}
\usepackage[T1]{fontenc}
\usepackage{lmodern,microtype,amsmath,amssymb,amsthm}
\usepackage{xcolor,fancyhdr}
\usepackage{xurl}
\usepackage[colorlinks=true,urlcolor=blue!45!black,linkcolor=blue!45!black]{hyperref}
\pagestyle{fancy}\fancyhf{}
\fancyhead[L]{Erd\H{o}s 708: the exact top-tail target}
\fancyhead[R]{23 November 2026}\fancyfoot[C]{\thepage}
\setlength{\headheight}{14pt}
\setlength{\emergencystretch}{2em}
\newcommand{\pos}[1]{(#1)_+}
\newcommand{\Top}{\operatorname{Top}}
\newcommand{\raw}{L_{\rm raw}}
\newcommand{\ind}{\mathbf1}
\newtheorem{proposition}{Proposition}[section]
\title{\vspace{-12mm}Erd\H{o}s 708: exact top tails\\and verified breaks in two intermediate arguments}
\author{Research checkpoint}
\date{23 November 2026}
\begin{document}
\maketitle\vspace{-5mm}
\noindent\textbf{Status.}
No uniform inequality proving \(g(n)\le(2+o(1))n\) was obtained.
No violation of that bound or of the exact \(2n\) bound was obtained.
The new results here are counterexamples to intermediate shortcuts,
with their scopes stated explicitly. Normalizations and exact
identities are not reported as a solution.

\section{The verified target and the exact remaining bank mass}
For \(X_w(t)=\sum_pw_pv_p(t)\), \(0\le w_p\le1\), write
\[
 H_I(w)=\sum_{b\in I}\pos{X_w(b)-1}.
\]
The required expression is exactly
\[
 Q_n(M,I,w)=X_w(M)+\Top_{n-1}\{X_w(a):2\le a<M\}-H_I(w).
 \tag{T}
\]
The forced term \(X_w(M)\) reflects the original condition
\(\max A=M\). Maximizing the fixed-set dual over all such source
sets gives (T), with no approximation.

Wang's raw finite-envelope theorem equates the worst-case leading
coefficients of the boxed raw LP and the integer problem
\cite{WangRaw}. Thus a uniform \(Q_n\le2n+o(n)\) would suffice.
His quadratic construction gives the exact \(2n\) bound for
\(M\ge16n^2\) \cite{WangMain}; its proof also works at fixed
ambient \(M\ge\max A\). No new rounding theorem is needed here.
The unresolved interval range is \(2n<M<16n^2\).

Only a \emph{jointly maximizing} source set and weight vector give
the paired optimality equations. Choosing a top set at an arbitrary
fixed weight vector does not supply those equations.
For a genuine paired optimum \((A,w,z)\), let
\[
 D_p(A)=\sum_{a\in A}v_p(a),\qquad
 V_p(U)=\sum_{b\in U}v_p(b),
\]
and let \(U\subseteq\{b\in I:X_w(b)>1\}\).
Every point of \(U\) has \(z_b=1\).
For each active prime, the whole raw row is tight, so
\[
 D_p(A)-V_p(U)
   =\sum_{b\in I\setminus U}z_bv_p(b)\ge0.
 \tag{1}
\]
Consequently the residual is exactly
\[
 \boxed{\ 
 \sum_pw_p\bigl(D_p(A)-V_p(U)\bigr)
 -\sum_{b\in I\setminus U}\pos{X_w(b)-1}
 =\sum_{b\in I\setminus U}z_b\ge0.\ }
 \tag{2}
\]
To verify (2), at a hot target \(z_b=1\), so the difference of its
two summands is one. At a target of load one the difference is
\(z_b\); at a cold target \(z_b=0\). These exhaust all positions.

Thus exact row balance does not discard the residual. It identifies
it as the still-selected mass outside the bank union. No uniform
bound of the required size for that mass was proved.
Similarly, a logarithmic cap must retain its unused source budget:
if \(E=F-\Phi_A(u)\) and \(S=(2+\delta)n-\Phi_A(u)\), then
\[
 E-S=F-(2+\delta)n.
 \tag{3}
\]
This equality is an accounting identity, not a bound on its sign.

\section{Dense same-threshold discrepancy is linear}
A possible dense shortcut would be to replace the target hinge by
the prefix hinge with an \(o(M)\) error. The next construction
disproves that shortcut, even with exact top-tail source selection
in the required numerical range.

Fix
\[
 \frac13<\alpha<\beta<\frac12,\qquad
 P_M=\{p\text{ prime}:M^\alpha<p\le M^\beta\},
 \qquad X(t)=\sum_{p\in P_M}v_p(t).
 \tag{4}
\]
All active weights are one. Put
\[
 k=|P_M|,\quad \mu_j=\sum_{p\in P_M}p^{-j},
 \quad \theta=\prod_{p\in P_M}(1-1/p).
 \tag{5}
\]
Every prefix integer has at most two active prime occurrences,
since \(3\alpha>1\). Therefore, with
\(N_2=|\{a\le M:X(a)=2\}|\), the exact formula is
\[
 H_{[1,M]}(w)=N_2
 =\sum_{\substack{p,q\in P_M\\p<q}}\left\lfloor\frac M{pq}\right\rfloor
  +\sum_{p\in P_M}\left\lfloor\frac M{p^2}\right\rfloor.
 \tag{6}
\]
In particular
\[
 \left|\frac{N_2}M-\frac{\mu_1^2+\mu_2}{2}\right|
 \le\frac{k(k+1)}{2M}=o(1).
 \tag{7}
\]
The last estimate follows already from \(k\le M^\beta\).
Mertens' product theorem \cite{Tao} gives
\[
 \mu_1\to\mu=\log(\beta/\alpha),\quad
 \theta\to e^{-\mu}=\alpha/\beta,\quad \mu_2,\mu_3\to0.
 \tag{8}
\]

\subsection*{A finite family of actual intervals}
For every \(p\in P_M\), choose independently and uniformly
\[
 r_p\in\{0,\ldots,p^3-1\},
\]
and let \(x\) be the least nonnegative solution of
\[
 x\equiv r_p\pmod{p^4}\qquad(p\in P_M).
 \tag{9}
\]
Use \(I_x=\{x+1,\ldots,x+M\}\). Since \(p^3>M\), every position
\(1\le j\le M\) satisfies \(1\le r_p+j<2p^3\le p^4\).
Its actual \(p\)-valuation is at most three, and for \(a=1,2,3\),
\[
 \mathbb P(p^a\mid x+j)=p^{-a}.
 \tag{10}
\]
These prime coordinates are independent at a fixed position.
No independence of different positions is used.

For integer-valued \(X\), \((X-1)_+=X-\ind_{X>0}\).
The exact finite average over these CRT choices is therefore
\[
 \boxed{\quad
 \mathbb E\,H_{I_x}(w)
 =M(\mu_1+\mu_2+\mu_3-1+\theta).
 \quad}
 \tag{11}
\]
Choose a member of this finite family whose hinge is at most its
average. Equations (6)--(11) imply
\[
 H_{[1,M]}(w)-H_{I_x}(w)
 \ge\left(\frac{\mu^2}{2}-\mu+1-e^{-\mu}+o(1)\right)M.
 \tag{12}
\]
\begin{samepage}
The coefficient
\[
 \kappa(\mu)=\frac{\mu^2}{2}-\mu+1-e^{-\mu}
 \tag{13}
\]
is strictly positive for \(\mu>0\): it vanishes at zero and has
derivative \(\mu-1+e^{-\mu}>0\).
\end{samepage}

For the concrete rational band
\[
 \alpha=\frac7{20},\quad\beta=\frac9{20},\quad\mu=\log(9/7),
\]
we have
\[
 \frac{\mu^2}{2}=0.03157947093\ldots,\qquad
 \kappa(\mu)=0.002487264872\ldots.
 \tag{14}
\]

\subsection*{The forced maximum and what is not disproved}
Let \(M\) tend to infinity through primes and set
\[
 A_M=\{a<M:X(a)=2\}\cup\{M\},\qquad n=N_2+1.
 \tag{15}
\]
Here \(X(M)=0\), and (15) is exactly the required top set with
the maximum forced in. Thus
\[
 Q_n(w)=2N_2-H_{I_x}(w)
 \ge\left(1+\frac{2\kappa(\mu)}{\mu^2}+o(1)\right)n
 =(1.07876208179\ldots+o(1))n.
 \tag{16}
\]
Weak duality makes this a lower bound for the raw covering optimum,
without claiming the chosen weights are optimal.
Also \(n/M\to0.03157947093\ldots\), and
\(2n<M<16n^2\) eventually.
The discrepancy (12) is therefore linear in \(n\) as well.
But
\[
 Q_n(w)\le2N_2<2n
\]
because the target hinge is nonnegative.
This construction refutes \(H_{[1,M]}-H_I=o(M)\) and the stronger
dense guess \(\raw\le n+o(n)\). It does \emph{not} refute
the desired coefficient two. Its sources have load two (and the forced maximum has load zero), so it does not meet the high-load minimal-obstruction hypothesis.

\section{A complete hot top tail still need not dominate its bank union}
This finite construction strengthens the earlier bank-product
counterexample to an exact top set, with every selected load above
two and the source maximum retained.

Let
\[
 N=11763=3^2\cdot1307,\quad h=5881,\quad p=11777,\quad
 M=pN=138532851.
 \tag{17}
\]
The supplied exact CRT certificate specifies residues \(c_r\) for
every prime \(r\le N\). Let \(C\) be the least nonnegative solution
of
\[
 C\equiv c_r\pmod r\quad(r\le N,\ r\text{ prime}),\qquad
 C\equiv h+1\pmod p.
 \tag{18}
\]
The full residue table and hexadecimal value of \(C\), together
with a standard-library verifier, are in the accompanying archive.
For \(J=\{C-h,\ldots,C+h\}\), that certificate checks:
\begin{itemize}
\item no member of \(J\) is divisible by \(p\);
\item exactly \(10352\) members have a prime divisor at most \(N/2\);
\item the other \(1411\) members avoid every prime at most \(N\);
\item the prefix \([4,N]\) has \(10353\) composite integers.
\end{itemize}
Set
\[
 I=\{p(C-h-1)+1,\ldots,p(C+h)\}.
 \tag{19}
\]
It has length \(M\); its \(p\)-multiples are exactly \(pt\),
\(t\in J\), and have \(p\)-valuation one.

Choose the exact clipped weights
\[
 w_p=1,\qquad w_r=\frac35\quad(r\le N,\ r\text{ prime}),
 \qquad w_r=0\ \text{otherwise},
 \tag{20}
\]
and take the \emph{whole} strict superlevel
\[
 A_*=\{a\in[2,M]:X_w(a)>2\}.
 \tag{21}
\]
Since \(p^2>M\), the selected \(p\)-multiples are exactly \(ps\)
with \(s\le N\) composite. Selected non-\(p\)-multiples have
at least four small-prime occurrences. Hence
\[
 \min_{a\in A_*}X_w(a)=\frac{11}{5},\qquad
 \max_{a\notin A_*}X_w(a)\le\frac95,\qquad
 X_w(M)=\frac{14}{5}.
 \tag{22}
\]
This is an exact top set, is closed upward for divisibility, and
retains every price-monotone exchange
\(a\mapsto ar/q\le M\) with \(q\mid a\), \(w_r\ge w_q\).

Let \(U\) be the union of \emph{all} banks \(a/q\), where
\(a\in A_*\) and \(q\mid a\) is prime. Then
\[
 \boxed{\quad D_p(A_*)=10353,\qquad V_p(U)=10352.\quad}
 \tag{23}
\]
Indeed, if \(t\in J\) is one of the \(1411\) rough positions,
\(X_w(pt)=1\). Every deleted source has
\(X_w(a/q)\ge11/5-1=6/5>1\), so no such quotient divides \(pt\).
Conversely, if a prime \(r\le N/2\) divides \(t\), the source
\(2rp\le M\) has load \(11/5\); deleting its occurrence of \(2\)
gives \(rp\mid pt\). This identifies the \(p\)-part of \(U\)
exactly and proves (23).

\subsection*{The numerical range and the exact failed optimality condition}
An integer prime-power sieve, supplied with its output, gives
\[
 n=|A_*|=56063090,\qquad D_2(A_*)=101157183.
 \tag{24}
\]
It counts this stated set and its row demand; it does not compute
\(g(n)\) or solve a covering problem. The mask is exactly
\(\Omega_{\{r\le N\}}(a)\ge4\), together with \(p\)-multiples
having at least two small-prime occurrences. Thus
\[
 2n=112126180<M=138532851<16n^2,\qquad n/M=0.4046916640\ldots.
 \tag{25}
\]

The displayed weights are \emph{provably nonoptimal}, not merely
unverified. Every multiple of \(4\) has load at least \(6/5\).
If these weights were optimal, all such targets would be selected
with \(z_b=1\). Their raw two-row supply is at least
\[
 2\left\lfloor\frac M4\right\rfloor+
 \sum_{j\ge3}\left\lfloor\frac M{2^j}\right\rfloor
 =103899623
 =D_2(A_*)+2742440.
 \tag{26}
\]
This contradicts tightness of the active row \(w_2=3/5>0\).
In fact the bank of \(4\) is itself in \(U\), from the source \(4p\).

Thus (23) refutes coordinatewise bank-product dominance even after
adding exact top selection, a strict load gap above two, divisibility
closure, all price-monotone exchanges, and the required numerical
range. Equation (26) shows exactly why it does not refute a statement
restricted to actual paired optima or to minimal obstructions.
The one-unit deficit in (23) is not an asymptotic counterexample.

\section{The unresolved step}
The complete top-tail and whole-row conditions were retained in
the attempted proof. They did not yield a uniform upper bound for
the residual selected mass in (2), either when \(n/M\to0\) or at
fixed positive density below the trivial interval-size threshold.
The dense same-threshold error cannot simply be declared negligible
by (12). The coordinatewise bank dominance shortcut fails by (23),
and its repair must use actual joint balance, as (26) demonstrates.

No new inequality sufficient for the required contradiction after
Wang's transfer is claimed. Exact \(2n\) remains separate.

\begin{thebibliography}{9}
\bibitem{WangRaw} Z. Wang, \emph{Raw fractional covers and worst-case
leading coefficients}, raw-transfer.tex in the interval-divisor-covers
repository, accessed 23 November 2026.
\url{https://github.com/science-prince/erdos-708-interval-divisor-covers/blob/main/papers/main-paper/raw-transfer.tex}.
\bibitem{WangMain} Z. Wang, \emph{Fractional-to-Integer Transference and
Five-Prime Bounds for Interval Divisor Covers}, Appendix on the
quadratic gate, accessed 23 November 2026.
\url{https://github.com/science-prince/erdos-708-interval-divisor-covers/blob/main/papers/main-paper/main.tex}.
\bibitem{Tao} T. Tao, \emph{Mertens' theorems}, 11 December 2013.
\url{https://terrytao.wordpress.com/2013/12/11/mertens-theorems/}.
\end{thebibliography}
\end{document}
